\documentclass{article}

% ================= PACKAGES =================
\usepackage[utf8]{inputenc}
\usepackage[T1]{fontenc}
\usepackage{amsmath, amssymb, bbm}
\usepackage{geometry, fancyhdr, hyperref, graphicx, float, booktabs, subcaption}
\usepackage{xcolor, ragged2e}
\geometry{a4paper, margin=1in}

% ================= HEADER & FOOTER =================
\pagestyle{fancy}
\fancyhf{}
\fancyhead[L]{\textbf{MATH 517 - Statistical Computation and Visualisation}}
\fancyhead[R]{Project Write-up}
\renewcommand{\headrulewidth}{0.4pt}
\fancyfoot[L]{} % Assuming this is the submitter, adjust if needed
\fancyfoot[R]{\textbf{Page \thepage}}
\renewcommand{\footrulewidth}{0.4pt}

\setlength{\parindent}{0pt}
\setlength{\parskip}{0.5em}

% ================= DOCUMENT =================
\begin{document}

\begin{titlepage}
\centering
\vspace*{4cm}
\rule{\textwidth}{1pt}\vspace{0.2cm}

{\LARGE\bfseries The EM Algorithm and Missing Data:\\[0.2cm] 
Performance Analysis and Comparison with Imputation Strategies}\\[0.3cm]
{\Large Students: Leonardo Tonelli, Luca Civilla, Valerio Viscovo}\\[0.2cm]
{\Large MATH 517 - Statistical Computation and Visualisation}\\[0.3cm]
{\large \today}\\[0.2cm]

\rule{\textwidth}{1pt}
\vfill
{\large \textbf{École Polytechnique Fédérale de Lausanne}}
\end{titlepage}

% ================= CONTENT =================
\section{Introduction}

Missing data is a ubiquitous problem in real-world statistical inference, often leading to biased estimates and reduced statistical power if handled via naive methods such as Complete Case Analysis (deleting rows with missing values). The Expectation-Maximization (EM) algorithm stands out as a powerful, iterative technique for finding Maximum Likelihood Estimates (MLE) of parameters in probabilistic models, particularly in the presence of unobserved latent variables or missing data.

\subsection{The EM Algorithm: Overview and Limitations}
The EM algorithm alternates between two steps: the \textbf{Expectation (E) step}, which computes the expected value of the log-likelihood function with respect to the conditional distribution of the missing data given the observed data and current parameter estimates; and the \textbf{Maximization (M) step}, which updates the parameters to maximize this expected log-likelihood.

While the EM algorithm is mathematically elegant and guaranties monotonic convergence of the likelihood, its practical deployment warrants careful consideration. Indeed, while EM is the gold-standard solution for MLE in the presence of missing data, faster and less complex imputation methods are often preferred for their ease of implementation and lower dependence on strict distributional assumptions. This project, therefore, addresses a critical objective: empirically demonstrating the operational boundary where EM's rigorous theoretical complexity yields an actual gain in accuracy over these more practical alternatives and assessing its robustness as core distributional and missing-data assumptions are stressed by varying degrees and patterns of missingness.

\section{Project Scope and Objectives}

The primary objective of this project is to answer the question: \textit{When faced with a missing data problem, under what conditions is the EM algorithm the optimal choice for statistical inference compared to imputation-based strategies?}

To answer this, we will analyze the performance of the EM algorithm across three dimensions: the mechanism of missingness, the proportion of missing data, and the underlying data distribution.

\subsection{Study Design: Synthetic and Real Data}
We will begin our analysis using a \textbf{synthetic Multivariate Gaussian dataset}. This controlled environment allows us to know the "ground truth" parameters ($\mu, \Sigma$), facilitating a precise calculation of estimation error.

However, real-world data rarely follows a perfect Gaussian distribution. To assess practical applicability, we will extend our analysis to real-world datasets. We propose using:

\begin{enumerate}
    \item \textbf{Fisher's Iris Dataset:} A classic dataset used for an initial low-dimensional test where the Gaussian assumption is most plausible. This allows for a baseline validation of the standard EM algorithm.
    \item \textbf{Wine Quality Dataset (UCI):}  A higher-dimensional dataset crucial for testing the performance of EM and imputation methods in a context characterized by non-linear correlations and potential multicollinearity. This complexity significantly stresses the assumptions of the standard Gaussian EM.
    \item \textbf{Boston Housing Dataset:} Included specifically to test performance on data that may contain heavier tails or outliers. This serves as a vital test for robustness in non-Gaussian contexts, preparing the ground for extensions using models like the t-Student EM.
\end{enumerate}

\section{Methodology}

\subsection{Missing Data Mechanisms and Proportions}
We will simulate missing values in our complete datasets (both synthetic and real) to evaluate performance against the ground truth. We will generate missingness masks based on Rubin’s three mechanisms:
\begin{itemize}
    \item \textbf{MCAR (Missing Completely At Random):} The probability of being missing is independent of both observed and unobserved data.
    \item \textbf{MAR (Missing At Random):} The probability of being missing depends on the observed data but not the missing data itself.
    \item \textbf{MNAR (Missing Not At Random):} The probability of being missing depends on the unobserved missing values. This is the most challenging scenario for standard inference.
\end{itemize}

For each mechanism, we will test varying degrees of severity by removing \textbf{10\%, 30\%, 50\%, and 75\%} of the data.

\subsection{Comparison Strategy: EM vs. Imputation}
A common alternative to the EM algorithm is to impute missing values to create a "complete" dataset, followed by standard Maximum Likelihood Estimation. We will compare the parameter estimates obtained via the direct EM algorithm against estimates obtained after applying the following imputation techniques:

\begin{itemize}
    \item \textbf{Baseline Methods:} Mean/Median/Mode imputation.
    \item \textbf{Regression Imputation:} Predicting missing values using observed covariates.
    \item \textbf{k-Nearest Neighbors (kNN):} Imputing based on local similarity in the feature space.
    \item \textbf{Multiple Imputation (MICE):} A stochastic approach that accounts for uncertainty in the missing values.
\end{itemize}

\subsection{Robustness Extension: The t-Student EM}
As an optional extension (time permitting), we investigate if the Gaussian EM is too rigid for real data. We plan to implement or utilize an EM algorithm based on the \textbf{multivariate t-student distribution}. We hypothesize that the t-distribution, with its heavier tails, may provide more robust parameter estimates in the presence of outliers or MNAR mechanisms compared to the standard Gaussian EM.

\section{Evaluation Metrics}
To quantitatively compare the EM algorithm with imputation methods, we will utilize the following metrics:
\begin{itemize}
    \item \textbf{Parameter Bias and MSE:} For synthetic data, we will calculate the Mean Squared Error between the estimated parameters ($\hat{\mu}, \hat{\Sigma}$) and the true parameters.
    \item \textbf{Log-Likelihood on Held-out Data:} For real data, where true parameters are unknown, we will evaluate the likelihood of a held-out test set given the parameters estimated by EM vs. Imputation.
    \item \textbf{Convergence Time:} Computational efficiency of EM versus iterative imputation methods like MICE.
\end{itemize}

By systematically varying the missingness patterns and comparing these approaches, we aim to provide a comprehensive guideline on when to deploy EM versus imputation for statistical inference.

\end{document}